% Generated shared reviewed material.
\section{Objective semantics, precision floors and replay}\label{ar:semantics}
The original true-circle sample-label least-squares objective is
\[
 \mathcal L_y(b)=\frac1{2m}\sum_{j=0}^{m-1}
 \left|\sum_{n=0}^d b_n\zeta_j^n-y_j\right|^2.
\]
With $d<m$, Fourier orthogonality gives $A^*A=mI$ and the real Euclidean
coordinate gradient $\nabla\mathcal L_y(b)=b-b^*$. The executed objective
$J$ instead has the exact rational midpoint target $t$.
The difference $\nabla\mathcal L_y-\nabla J=t-b^*$ is enclosed
coefficientwise by $e_n$.
Consequently this is an actual rational GD run for $J$ and a certified
perturbed-gradient implementation for $\mathcal L_y$. It is not the exact
optimizer trajectory of the ideal-root objective, nor a replay of neural or
quantum training. An interval enclosure of an ideal-root GD trajectory
would be a different explicitly labelled object.

Each actual update contracts $b^k-t$ by $3/4$, so the exact stored losses
satisfy $J(b^{k+1})=(9/16)J(b^k)$. A replay recomputes every rational
coordinate update, loss and weighted residual. The closed form
$b^k=(1-(3/4)^k)t$ is an independent check on the executed trace, not its
producer. Exact rational loop coordinates introduce no further midpoint or
floating export error. The arithmetic budget covers the difference between
the certified Fourier rectangles and their midpoint targets, including the
inherited root enclosure and outward interval operations. A subsequent
float conversion requires its own exact conversion-error bound.

The normalized polynomial convention is essential: a stored coefficient
with index $n$ multiplies $((s-1/2)/(6/5))^n$, not $(s-1/2)^n$.
The input certificate is pinned by SHA-256, including its source hashes,
all $192$ sample enclosures, root rectangles and analytic radius components.
Producing source is captured at entry, checked unchanged before emission,
and copied into a fresh output directory. An existing output is never
overwritten. An independent replay should also check the fixed input hash,
all dimensions, degree choice, the first stopping iterate, all five exact
terms and their literal sum against the requested rational tolerance.

The attainable boundary of this majorant is
\[
 C_{\min,q}=\min_{0\le d<64}C_{q,d}.
\]
Every request $\varepsilon>C_{\min,q}$ has strict budget slack for some
degree and the proved finite stop rule. At or below that boundary this
interface refuses. The degree minimizing $C_d$ need not be $63$:
Taylor error decreases with degree while the oracle and arithmetic sums
increase. In particular $C_d\ge\eta$ and
$C_d\ge M_q2^{-64}/(1-2^{-64})$. Increasing the number of optimization
steps does not lower these fixed-data budgets. Approximate boundary values
on the stated compact are $7.2929\cdot10^{-17}$,
$1.7370\cdot10^{-16}$ and $3.3741\cdot10^{-16}$ for $q=3,4,5$;
only the exact exported rational minima are used for acceptance.
These are certificate floors, not lower bounds on the actual errors.

A full adaptive theta oracle would separately need proof-preserving choices
of omitted-$n$ cutoff, time cutoff, quadrature order, whole-panel discard
budget, primitive series lengths, interval precision and serialized-bound
precision. The present $N=11$, $T=128$, $224$-bit arithmetic,
$60/92/60$ primitive series lengths, $2^{-90}$ discard threshold and
$180$-bit rounded budgets cannot support every positive request.
Increasing a nominal bit count alone leaves the fixed primitive remainders
and analytic tails in place. The preserved labels are not recomputed here.
The classical identity, finite arithmetic and informal proof are not new
Lean declarations. No zero-location, all-conductor, positive arithmetic
determinant, quantum-advantage or GRH conclusion follows.

The interactive replay is available at \href{https://amelie-iska.github.io/GRH-research/readout.html}{the certified readout laboratory}. All decimal graphics are illustrative views of the linked exact records.
